\subsection{序列空间中的对角算子}

设 I 是非空集合。回顾（EVT, I, p. 4 与 p. 5），由有界族 \(x = (x_i)_{i \in \mathrm{I}}\)（\(\mathbf{C}\) 的元素）构成的向量空间，赋以范数 \(\| x \| = \sup_i |x_i|\)，是记作 \(\ell_{\mathbf{C}}^{\infty}(\mathrm{I})\) 的巴拿赫空间。

设 \(p\) 是 \([1, +\infty[\) 的一个元素。空间 \(\ell_{\mathbf{C}}^{p}(\mathrm{I})\)（由族 \(x = (x_{i})_{i \in \mathrm{I}}\)（\(\mathbf{C}\) 的元素，且使族 \((|x_{i}|^{p})\) 可和）构成）是一个向量空间，在其上映射 \(x \mapsto \| x \| = (\sum_{i} |x_{i}|^{p})^{1/p}\) 是一个范数，且使它成为巴拿赫空间。若给集合 I 赋以离散拓扑以及使 \(\mu(\{x\}) = 1\)（对一切 \(x \in \mathrm{I}\)）的测度，这个空间正是巴拿赫空间 \(\mathrm{L}_{\mathbf{C}}^{p}(\mathrm{I}, \mu)\)。当 \(p = 1\) 或 p = 2 时，这个记号于是与 EVT, I, p. 4 及 EVT, V, p. 18 中的记号一致（参看 INT, IV, p. 141, §4, 第 1 目, 例）。

对复数的族 \(x = (x_{i})_{i \in \mathrm{I}}\) 与 \(y = (y_{i})_{i \in \mathrm{I}}\)，用 xy 表示族 \((x_{i}y_{i})_{i \in \mathrm{I}}\)。

命题 1. --- 设 \(\lambda = (\lambda_i)_{i\in I}\) 是 C 的元素的一个有界族。设 \(p\) 是 \([1, +\infty]\) 的一个元素。设 E 是巴拿赫空间 \(\ell_{\mathbf{C}}^p (\mathrm{I})\)。

\begin{enumerate}
\def\labelenumi{\alph{enumi})}
\item
  对每个 \(x \in E\)，有 \(\lambda x \in E\)，且映射 u：\(x \mapsto \lambda x\) 是 E 的连续自同态，它在巴拿赫代数 \(\mathcal{L}(E)\) 中的谱等于诸 \(\lambda_{i}\) 所成集合在 C 中的闭包，其范数等于 \(\sup_{i} |\lambda_{i}|\)；
\item
  自同态 u 是紧的，当且仅当对每个实数 \(\varepsilon > 0\)，元素 \(i \in I\) 中满足 \(|\lambda_{i}| \geqslant \varepsilon\) 者所成的集合是有限的。
\end{enumerate}

采用命题中的记号，自同态 \(x \mapsto \lambda x\) 称为 \(\ell_{\mathbf{C}}^{p}(\mathrm{I})\) 的对角自同态。

我们来证明该命题。设 \(C = \sup_{i} |\lambda_{i}|\)。设 \(x \in E\)。我们有不等式

\[
\sum_ {i \in I} | \lambda_ {i} x _ {i} | ^ {p} \leqslant C ^ {p} \| x \| ^ {p}, \text {   if   } p \neq + \infty
\]

\[
\sup _ {i \in I} | \lambda_ {i} x _ {i} | \leqslant C \| x \|, \text {   if   } p = + \infty .
\]

这就证明了 \(\lambda x \in E\)。于是映射 \(x \mapsto \lambda x\) 是 E 的自同态，且这些不等式证明它具有范数 \(\leqslant C\)。

对 \(j \in \mathrm{I}\)，用 \(e_j\) 表示元素 \((x_i)_{i \in \mathrm{I}}\)（\(\ell_{\mathbf{C}}^p(\mathrm{I})\) 中的）：它满足 \(x_j = 1\)，且 \(x_i = 0\)（当 \(i \neq j\)）。于是有 \(u(e_j) = \lambda_j e_j\)（对一切 \(j \in \mathrm{I}\)），这表明 \(\lambda_j\) 属于 \(u\) 的谱。由于 \(u\) 的谱是闭的，在 \(\mathbf{C}\) 中诸 \(\lambda_i\) 所成集合的闭包包含于 \(u\) 的谱。由于 \(u\) 的谱包含在以 0 为中心、以 \(\|u\|\) 为半径的圆盘中（I, p. 24, 定理 1 与 I, p. 21 的公式 (3)），我们推出不等式 \(C \leqslant \|u\|\)，从而 \(\|u\| = C\)。

反之，若 \(\lambda \in \mathbf{C}\) 不附着于诸 \(\lambda_{i}\) 所成的集合，则族 \(((\lambda_{i} - \lambda)^{-1})_{i\in \mathrm{I}}\) 有界。于是前面的论证表明，线性映射 \(x\mapsto ((\lambda -\lambda_i)^{-1}x_i)_{i\in \mathrm{I}}\) 是 E 的一个自同态。它是 \(u - \lambda 1_{\mathrm{E}}\) 的逆，因而 \(\lambda\) 不属于 \(u\) 的谱。这就证明了 \(a\)）。

现在证明 \(b\)）。假设自同态 \(u\) 是紧的。设 \(\varepsilon > 0\)。用 J 表示 \(i \in I\) 中满足 \(|\lambda_i| \geqslant \varepsilon\) 者所成的集合。

诸元素 \(e_i\)（\(i \in \mathbf{J}\)）所成的集合在 E 中有界。它在 \(u\) 下的像由诸元素 \(\lambda_i e_i\)（\(i \in \mathbf{J}\)）组成，因此在 E 中相对紧（III, p. 2 的注记 1）。由于 \(\| \lambda_i e_i - \lambda_j e_j \| \geqslant \varepsilon\)（对 J 中每一对元素 \(i \neq j\)），这仅当集合 J 有限时才可能。

我们来证明逆命题。设 J 是 I 的一个有限子集，记 \(z_{\mathrm{J}} = (z_{\mathrm{J},i})_{i \in \mathrm{I}}\)，其中 \(z_{J,i} = \lambda_i\)（当 \(i \in J\)），而 \(z_{J,i} = 0\)（当 \(i \in I - J\)）；族 \(z_J\) 定义了 E 的一个有限秩自同态 \(u_J \colon x \mapsto z_J x\)。由前所证，线性映射 \(u - u_J\) 的范数被 \(R_+\) 中族 \((|\lambda_i|)_{i \in I-J}\) 的上确界所控制。该假设蕴含：对每个 \(\varepsilon > 0\)，存在有限子集 \(J \subset I\)，使得这个上界 \(\leqslant \varepsilon\)。由此推出映射 u 附着于 \(\mathcal{L}^{\mathrm{f}}(\mathrm{E})\)，因而它是紧的（III, p. 4, 命题 2 的推论）。

注记. --- 1) 若 I 有限，断言 b) 的假设总成立。当集合 I 无穷时，该假设意味着族 \((\lambda_{i})_{i\in\mathrm{I}}\) 沿 I 的有限子集的余集所成的滤子趋于 0；特别地，\(i\in I\) 中满足 \(\lambda_{i}\neq0\) 者所成的集合这时是可数的。

\begin{enumerate}
\def\labelenumi{\arabic{enumi})}
\setcounter{enumi}{1}
\tightlist
\item
  我们稍后将看到（参看 IV, p. 149, 定理 1），当 \(p = 2\) 时，希尔伯特型空间 \(\ell_{\mathbf{C}}^{2}(\mathrm{I})\) 的每个紧正规自同态都形如 \(u = w \circ v \circ w^{-1}\)，其中 \(v\) 是紧对角自同态，而 \(w\) 是希尔伯特空间 \(\ell_{\mathbf{C}}^{2}(\mathrm{I}).*\) 的酉自同态（EVT, V, p. 40）。
\end{enumerate}

